Large language models(LLMs)는 data-processing workflow 자동화에 점점 더 많이 사용되고 있지만, coding agent가 생성한 script는 일반적으로 persistent하고 editable한 platform artifact로 자동 materialize되지 않는다. 우리는 이 disconnect를 \textit{NL2Pipeline gap}이라고 부른다. 이를 해소하기 위해 \textsc{DataFlow-Harness}를 제안한다. 이 platform은 free-form script 대신 typed, incremental mutation을 통해 LLM agent가 platform-native directed acyclic graph(DAG)를 구성하도록 안내한다. Platform은 procedural guidance를 제공하는 \textsc{DataFlow-Skills}, live operator registry와 현재 pipeline state를 노출하는 Model Context Protocol(MCP) layer, conversational authoring과 visual DAG editor를 동기화하는 \textsc{DataFlow-WebUI}를 결합한다. 12개 task로 구성된 data-engineering benchmark에서 \textsc{DataFlow-Harness}는 관찰된 end-to-end pass rate 93.3\%를 달성한다. Vanilla Claude Code와 비교하면 측정된 monetary cost를 72.5\%, generation latency를 49.9\% 줄인다. Context-Aware Claude Code baseline과 비교하면 관찰된 pass rate는 0.9 percentage point 이내이며, cost는 42.8\% 더 낮다. Task별 분석은 construction이 암묵적 procedural knowledge에 의존할 때 Skills가 가장 유용함을 보여준다. 이러한 결과는 live platform grounding이 script-generation baseline에 가까운 관찰 reliability를 유지하면서도 더 낮은 measured construction cost와 latency로 persistent하고 editable한 workflow artifact를 생성할 수 있음을 보여준다.
